% ASSERT_CONVENTION: natural_units=natural, metric_signature=mostly_minus, fourier_convention=physics, coupling_convention=alpha_s, generator_normalization=T(fund)=1/2, covariant_derivative_sign=D_mu=partial_mu-igA
%
% Exercise Set 5: Fermion Mass Hierarchy from Scalar Democracy
% Part of: The Dualised Standard Model -- Part III Exercises
%
% Conventions (Seiberg hep-th/9411149, unchanged from Lectures 1-8):
%   T(fund) = 1/2 per Weyl fermion
%   b_0 = (11N_c - 2N_f)/3 for N_f Dirac flavours
%   Universal Yukawa: y (single coupling for all fermions)
%   VEV convention: v^2 = sum_{ij} (|H^u_{ij}|^2 + |H^d_{ij}|^2) = (246 GeV)^2/2
%   Epistemic: all non-SUSY duality claims carry conjectural qualifiers (P1)
%
\documentclass[11pt,a4paper]{article}
% ASSERT_CONVENTION: natural_units=natural, metric_signature=mostly_minus, fourier_convention=physics, coupling_convention=alpha_s, generator_normalization=T(fund)=1/2, covariant_derivative_sign=D_mu=partial_mu-igA
%
% CONVENTIONS (Seiberg hep-th/9411149)
% Metric: (+,-,-,-) (mostly minus)
% T(fund) = 1/2 per Weyl fermion
% b_0 = (11N_c - 2N_f)/3 for N_f Dirac flavours
% D_mu = partial_mu - ig A_mu^a T^a
% A(fund) = 1 (cubic anomaly coefficient per fundamental Weyl fermion)
% alpha_s = g_s^2 / (4 pi)
% R[Q] = (N_f - N_c)/N_f; R[W] = 2
% N_f counts Dirac flavour pairs (each = 2 Weyl fermions)
%
% This file is \input{} by all lecture files.

%% ---------- Document class ----------
% (loaded by the wrapper; preamble only provides packages and macros)

%% ---------- Standard packages ----------
\usepackage{amsmath,amssymb,amsthm,mathtools}
\usepackage{hyperref}
\usepackage[capitalise,noabbrev]{cleveref}
\usepackage{enumitem}
\usepackage{booktabs}
\usepackage{array}
\usepackage{xcolor}
\usepackage{tikz}
\usetikzlibrary{arrows.meta,decorations.markings,positioning,calc}
\usepackage{fancybox}
\usepackage{tcolorbox}
\tcbuselibrary{skins,breakable}
\usepackage{slashed}              % Feynman slash notation
\usepackage{cite}                 % compressed citation lists

%% ---------- Hyperref configuration ----------
\hypersetup{
  colorlinks=true,
  linkcolor=blue!70!black,
  citecolor=green!50!black,
  urlcolor=blue!80!black,
  pdfauthor={Part III Lecture Notes},
  pdftitle={The Dualised Standard Model}
}

%% ---------- Theorem environments ----------
\theoremstyle{plain}
\newtheorem{theorem}{Theorem}[section]
\newtheorem{lemma}[theorem]{Lemma}
\newtheorem{proposition}[theorem]{Proposition}
\newtheorem{corollary}[theorem]{Corollary}

\theoremstyle{definition}
\newtheorem{definition}[theorem]{Definition}
\newtheorem{example}[theorem]{Example}
\newtheorem{exercise}{Exercise}[section]

\theoremstyle{remark}
\newtheorem{remark}[theorem]{Remark}

%% ---------- Physics macros: gauge groups and representations ----------
\newcommand{\Nc}{N_{\!c}}
\newcommand{\Nf}{N_{\!f}}
\newcommand{\SU}[1]{\mathrm{SU}(#1)}
\newcommand{\UU}[1]{\mathrm{U}(#1)}

% Representations
\newcommand{\fund}{\square}
\newcommand{\antifund}{\overline{\square}}
\newcommand{\adj}{\mathrm{adj}}
\newcommand{\antisym}{\wedge^{\!2}}        % antisymmetric rank-2
\newcommand{\sym}{\mathrm{S}_2}           % symmetric rank-2

%% ---------- Physics macros: group theory constants ----------
% Dynkin index: Tr(T^a T^b) = T(R) delta^{ab}
\newcommand{\Tfund}{\tfrac{1}{2}}          % T(fund) = 1/2 per Weyl fermion
\newcommand{\Tadj}{\Nc}                    % T(adj) = N_c
\newcommand{\Cfund}{\frac{\Nc^2 - 1}{2\Nc}}  % C_2(fund)
\newcommand{\Cadj}{\Nc}                    % C_2(adj)

% Cubic anomaly coefficient: A(R) = Tr_R[T^a {T^b, T^c}] with A(fund) = 1
\newcommand{\Afund}{1}

%% ---------- Physics macros: beta function ----------
\newcommand{\bzero}{b_0}
\newcommand{\bone}{b_1}
\newcommand{\alphas}{\alpha_s}
\newcommand{\gs}{g_s}

%% ---------- Physics macros: SUSY / R-charge ----------
\newcommand{\Rcharge}[1]{R[#1]}
\newcommand{\Wsup}{W}                      % superpotential

%% ---------- Physics macros: covariant derivative and field strength ----------
\newcommand{\Dmu}{D_\mu}
\newcommand{\Fmn}{F_{\mu\nu}}
\newcommand{\Ftilde}{\widetilde{F}_{\mu\nu}}

%% ---------- Physics macros: common abbreviations ----------
\newcommand{\ie}{\textit{i.e.}}
\newcommand{\eg}{\textit{e.g.}}
\newcommand{\cf}{\textit{cf.}}
\newcommand{\IR}{\ensuremath{\mathrm{IR}}}
\newcommand{\UV}{\ensuremath{\mathrm{UV}}}
\newcommand{\AF}{\ensuremath{\mathrm{AF}}}
\newcommand{\BZ}{\ensuremath{\mathrm{BZ}}}
\newcommand{\NSVZ}{\ensuremath{\mathrm{NSVZ}}}
\newcommand{\SQCD}{\ensuremath{\mathrm{SQCD}}}
\newcommand{\QCD}{\ensuremath{\mathrm{QCD}}}
\newcommand{\SM}{\ensuremath{\mathrm{SM}}}
\newcommand{\Tr}{\mathrm{Tr}}

%% ---------- Convention box macro ----------
\newtcolorbox{conventionbox}[1][]{%
  enhanced,
  colback=blue!5!white,
  colframe=blue!60!black,
  fonttitle=\bfseries,
  title={Conventions},
  #1
}

%% ---------- Comparison table style ----------
\newtcolorbox{comparisontable}[1][]{%
  enhanced,
  colback=orange!5!white,
  colframe=orange!70!black,
  fonttitle=\bfseries,
  title={Comparison},
  #1
}

%% ---------- Warning / important box ----------
\newtcolorbox{warningbox}[1][]{%
  enhanced,
  colback=red!5!white,
  colframe=red!70!black,
  fonttitle=\bfseries,
  title={Important},
  #1
}

%% ---------- Preview / forward-reference box ----------
\newtcolorbox{previewbox}[1][]{%
  enhanced,
  colback=green!5!white,
  colframe=green!50!black,
  fonttitle=\bfseries,
  title={Preview},
  #1
}

%% ---------- Page layout ----------
\usepackage[margin=2.5cm]{geometry}
\setlength{\parskip}{0.4em}
\setlength{\parindent}{0em}

%% ---------- Section formatting ----------
\usepackage{titlesec}
\titleformat{\section}{\Large\bfseries}{\thesection}{1em}{}
\titleformat{\subsection}{\large\bfseries}{\thesubsection}{1em}{}

%% ---------- End of preamble ----------


\title{\textbf{Exercise Set 5: Fermion Mass Hierarchy from Scalar Democracy}}
\author{The Dualised Standard Model --- Part III Exercises}
\date{Lent Term 2027}

\begin{document}
\maketitle

\begin{conventionbox}[title={Conventions for these exercises}]
\renewcommand{\arraystretch}{1.3}
\begin{tabular}{@{}l l@{}}
\toprule
\textbf{Quantity} & \textbf{Convention} \\
\midrule
Units & $\hbar = c = 1$ (natural) \\
Metric & $\eta_{\mu\nu} = \mathrm{diag}(+1,-1,-1,-1)$ (mostly minus) \\
Universal Yukawa & $y$ (single magnetic coupling for all fermions) \\
VEV sum rule & $v^2 = \sum_{i,j}\bigl(\langle H_{ij}^u\rangle^2
  + \langle H_{ij}^d\rangle^2\bigr) = \tfrac{1}{2}(246\;\text{GeV})^2$ \\
Planck mass & $M_{\text{Pl}} = 1.22 \times 10^{19}\;\text{GeV}$ (reduced: $M_P = M_{\text{Pl}}/\sqrt{8\pi}$) \\
Epistemic & ``if the conjectured duality holds'' for all non-SUSY claims (P1) \\
\bottomrule
\end{tabular}

\medskip

\noindent References: Lecture~7 (\S\S2--3); Lecture~8 (\S\S1--2);
Cacciapaglia--Sannino~\cite{CacciapagliaSannino2024};
Hill, Machado, Thomsen, Turner~\cite{HillMachado2019}.

\medskip

\noindent\textbf{Epistemic note.}
All results in this exercise assume the conjectured non-SUSY
duality holds.  The group-theoretic statements (VEV sum rule,
effective Yukawa structure) are exact consequences of the magnetic
Lagrangian.  The physical interpretation as a fermion mass hierarchy
mechanism is conjectural.
\end{conventionbox}

\bigskip


%% ========================================================================
%%  PROBLEM 1: Top Quark Mass from Universal Yukawa Coupling
%% ========================================================================
\section*{Problem 1: Top Quark Mass from Universal Yukawa
  \hfill $\star\star$ \normalsize(15 min, 25 marks)}
\addcontentsline{toc}{section}{Problem 1: Top quark mass from universal Yukawa}
\label{prob:top-mass}

\noindent\textbf{Background.}
In the dualised Standard Model (Lectures~7--8), \emph{if the conjectured
duality holds}, the magnetic Yukawa interaction between quarks and the
18 Higgs doublets $H_{ij}^{u,d}$ ($i,j = 1,2,3$ labelling generations)
takes the form
\begin{equation}\label{eq:yukawa-universal}
  \mathcal{L}_Y
  \;=\; y \sum_{i,j=1}^{3}
  \Bigl(
    \bar{q}_L^{\,i}\, u_R^{\,j}\, H_{ij}^u
    \;+\;
    \bar{q}_L^{\,i}\, d_R^{\,j}\, H_{ij}^d
  \Bigr)
  \;+\; \text{h.c.}\,,
\end{equation}
where $y$ is a \textbf{single universal Yukawa coupling}, the same for
every generation pair.  This is a radical departure from the Standard
Model, where each fermion species has its own independent Yukawa
coupling.

The effective Yukawa couplings that determine fermion masses are
(Lecture~8, Eq.~(2)):
\begin{equation}\label{eq:eff-yukawa}
  Y_{ij}^u \;=\; y\, \frac{\langle H_{ij}^u \rangle}{v}\,,
  \qquad
  Y_{ij}^d \;=\; y\, \frac{\langle H_{ij}^d \rangle}{v}\,,
\end{equation}
where the total electroweak VEV satisfies
$v^2 = \sum_{ij}(\langle H_{ij}^u\rangle^2
+ \langle H_{ij}^d\rangle^2) = (246\;\text{GeV})^2/2$.

Planck-scale four-fermion operators (Lecture~8, \S2) generate mixing
between the leading Higgs doublet $H_0 \equiv H_{33}^u$ and the 17
dormant doublets $H_i$ ($i = 1, \ldots, 17$):
\begin{equation}\label{eq:planck-op}
  \mathcal{L}_{\text{Planck}}
  \;\supset\;
  \frac{c^{abcd}}{M_{\text{Pl}}^2}\,
  (Q^a \widetilde{Q}^b)(Q^c \widetilde{Q}^d)^\dagger\,,
\end{equation}
where $c^{abcd}$ are $\mathcal{O}(1)$ dimensionless coefficients.

\medskip

\noindent\textbf{Questions.}
\begin{enumerate}[label=(\alph*)]

  \item \textbf{[5 marks]}
    Suppose the leading VEV satisfies
    $\langle H_0 \rangle = \langle H_{33}^u \rangle \approx v$, with all
    other VEVs suppressed.  Show that the top quark Yukawa coupling is
    \begin{equation}\label{eq:top-yukawa-prob}
      y_t \;=\; y \,\frac{\langle H_{33}^u \rangle}{v}
      \;\approx\; y\,,
    \end{equation}
    and hence
    \begin{equation}\label{eq:top-mass-prob}
      m_t \;=\; y_t\, v \;\approx\; y\, v\,.
    \end{equation}
    For $y \sim \mathcal{O}(1)$ and $v \approx 174\;\text{GeV}$, show
    that $m_t \sim \mathcal{O}(v) \sim 174\;\text{GeV}$.  State this as
    a \textbf{parametric consistency check}: the observed
    $m_t = 172.7 \pm 0.3\;\text{GeV}$ requires $y \approx 0.99$, which
    is naturally $\mathcal{O}(1)$.  Note that $y$ is a \textbf{free
    parameter} of the magnetic theory.

  \item \textbf{[8 marks]}
    The dormant Higgs doublets $H_i$ have large positive masses
    $M_i^2 \gg v^2$ and acquire induced VEVs through mixing with the
    leading Higgs $H_0$ (Lecture~8, Eq.~(8)):
    \begin{equation}\label{eq:induced-vev-prob}
      \langle H_i \rangle
      \;=\; \frac{\mu_{0i}^2}{M_i^2}\, \langle H_0 \rangle\,,
    \end{equation}
    where $\mu_{0i}^2$ are off-diagonal mixing mass terms generated by
    the Planck operators.
    \begin{enumerate}[label=(\roman*)]
      \item Show that the effective Yukawa coupling for a fermion species
        $f$ whose mass arises from the Higgs doublet $H_i$ is:
        \begin{equation}\label{eq:eff-yukawa-dormant}
          y_f \;=\; y\, \frac{\mu_{0i}^2}{M_i^2}\,.
        \end{equation}

      \item From dimensional analysis of the Planck operators
        (Eq.~\eqref{eq:planck-op}), the mixing mass and dormant Higgs
        mass scale as
        \[
          \mu_{0i}^2 \;\sim\; \xi_i\, \frac{\Lambda_D^4}{M_{\text{Pl}}^2}\,,
          \qquad
          M_i^2 \;\sim\; \Lambda_D^2\,,
        \]
        where $\xi_i$ is an $\mathcal{O}(1)$ coefficient.  Show that
        \begin{equation}\label{eq:mass-ratio-suppression}
          \frac{m_f}{m_t}
          \;\sim\; \xi_i\, \frac{\Lambda_D^2}{M_{\text{Pl}}^2}\,.
        \end{equation}

      \item For $\Lambda_D = 10^{11}\;\text{GeV}$ and
        $M_{\text{Pl}} = 1.22 \times 10^{19}\;\text{GeV}$, evaluate
        $\Lambda_D^2 / M_{\text{Pl}}^2$ numerically.  Compare this
        suppression factor with the observed mass ratios $m_b/m_t
        \approx 0.024$, $m_c/m_t \approx 0.007$, and $m_s/m_t
        \approx 5 \times 10^{-4}$.  Comment on whether the parametric
        suppression is in the right ballpark.
    \end{enumerate}

  \item \textbf{[7 marks]}
    The bottom quark mass ratio $m_b/m_t \approx 0.024$ arises from the
    down-type Higgs VEV (Lecture~8, Eq.~(6)):
    \begin{equation}\label{eq:bottom-ratio-prob}
      \frac{m_b}{m_t}
      \;=\; \frac{\langle H_{33}^d \rangle}{\langle H_{33}^u \rangle}\,.
    \end{equation}

    \begin{enumerate}[label=(\roman*)]
      \item As shown in part~(b), the generic Planck suppression
        $(\Lambda_D/M_{\text{Pl}})^2 \sim 10^{-16}$ is far too small
        to produce $m_b/m_t \approx 0.024$.  Explain why the bottom
        quark mass instead arises from \textbf{direct}
        $H_{33}^u$--$H_{33}^d$ mixing through the scalar potential,
        with the suppression controlled by the ratio
        $\mu_b^2/M_{d,33}^2 \sim \mathcal{O}(10^{-2})$ of
        $\mathcal{O}(1)$ parameters in the Planck-scale operator
        coefficients.

      \item For the \emph{lighter} fermion generations (charm, strange,
        and below), the mass ratios arise from a cascade of mixing
        steps in the scalar potential.  Each step is controlled by a
        ratio $\mu_i^2/M_i^2$ of $\mathcal{O}(1)$ parameters.  Explain
        why the charm quark ($m_c/m_t \approx 0.007$) sits at a
        comparable suppression level to the bottom quark (i.e., the
        same ``tier'' of the VEV hierarchy), while the strange quark
        ($m_s/m_t \approx 5 \times 10^{-4}$) requires an additional
        step of suppression (a deeper tier).

      \item Assign the bottom, charm, and strange quarks to appropriate
        ``tiers'' in the VEV hierarchy.  Which quarks share the same
        tier, and which require additional steps of cascaded mixing?
    \end{enumerate}

  \item \textbf{[5 marks]}
    \textbf{Discussion.}
    \begin{enumerate}[label=(\roman*)]
      \item Is the fermion mass hierarchy a \textbf{prediction} of the
        dualised SM, or a \textbf{parametric consistency check}?  Justify
        your answer by listing all free parameters in the framework.

      \item What does the dualised SM \emph{explain} about the fermion
        mass hierarchy that the Standard Model does not?

      \item Identify the \textbf{weakest assumption} in the above
        derivation: which step would invalidate the entire argument if it
        failed?
    \end{enumerate}

\end{enumerate}

\medskip

\noindent\textit{Hints:}
\begin{itemize}
  \item For (a): remember that $v = 246/\sqrt{2}\;\text{GeV}
    \approx 174\;\text{GeV}$ when $v$ is defined as the VEV rather
    than $\sqrt{2}\, G_F^{-1/2}$.  Adapt to whichever convention is
    stated.
  \item For (b)(iii): compute $\log_{10}(\Lambda_D^2/M_{\text{Pl}}^2)$
    and compare with $\log_{10}(m_f/m_t)$ for each fermion.
  \item For (c)(ii): compare the suppression ratios $\mu_f^2/M_f^2$
    needed for each fermion and group them into tiers.
\end{itemize}

\bigskip
\hrule
\bigskip


%% ========================================================================
%%  SOLUTION TO PROBLEM 1
%% ========================================================================
\begin{proof}[\textbf{Solution to Problem 1}]

\textbf{(a) Top quark mass from universal Yukawa.}

From the effective Yukawa structure~\eqref{eq:eff-yukawa}, the top
quark Yukawa coupling is:
\[
  y_t \;=\; y\, \frac{\langle H_{33}^u \rangle}{v}\,.
\]
If the leading VEV dominates, $\langle H_{33}^u \rangle \approx v$,
then:
\[
  y_t \;\approx\; y \cdot \frac{v}{v} \;=\; y\,.
\]
The top quark mass is:
\[
  m_t \;=\; y_t\, v \;\approx\; y\, v\,.
\]

\textbf{Dimensional check:} $[y] = [\text{dimensionless}]$,
$[v] = [\text{GeV}]$, so $[m_t] = [\text{GeV}]$.
\quad$\checkmark$

\textbf{Numerical evaluation.}  With $v \approx 174\;\text{GeV}$:
\[
  m_t \;\approx\; y \times 174\;\text{GeV}\,.
\]
The measured top mass $m_t = 172.7 \pm 0.3\;\text{GeV}$ requires:
\[
  y \;\approx\; \frac{172.7}{174} \;\approx\; 0.99\,,
\]
which is naturally $\mathcal{O}(1)$.
\[
  \boxed{m_t \;\approx\; y\, v \;\sim\; \mathcal{O}(v)
  \;\sim\; 174\;\text{GeV}\,,
  \quad\text{with } y \approx 0.99 \sim \mathcal{O}(1)\,.}
\]

\textbf{Parametric consistency statement.}
This is \emph{not} a prediction: the universal Yukawa coupling $y$ is a
\textbf{free parameter} of the magnetic theory, determined by the
strong dynamics at the duality scale $\Lambda_D$.  The result
$m_t \sim \mathcal{O}(v)$ shows that $y_t \sim \mathcal{O}(1)$ is
\emph{natural} in the dualised SM---in contrast to the SM, where
$y_t \approx 1$ is an unexplained numerical coincidence.

\bigskip

\textbf{(b) Induced VEVs and mass suppression.}

\textit{(i) Effective Yukawa from dormant Higgs.}

A fermion species $f$ whose mass arises from the dormant Higgs $H_i$
has:
\[
  m_f \;=\; y\, \frac{\langle H_i \rangle}{v}\, v
  \;=\; y\, \langle H_i \rangle\,.
\]
Using the induced VEV~\eqref{eq:induced-vev-prob}:
\[
  m_f \;=\; y\, \frac{\mu_{0i}^2}{M_i^2}\, \langle H_0 \rangle
  \;=\; y\, \frac{\mu_{0i}^2}{M_i^2}\, v\,.
\]
Hence the effective Yukawa coupling is:
\[
  y_f \;\equiv\; \frac{m_f}{v}
  \;=\; y\, \frac{\mu_{0i}^2}{M_i^2}\,.
\]
\[
  \boxed{y_f \;=\; y\, \frac{\mu_{0i}^2}{M_i^2}\,.}
  \quad\checkmark
\]

\textit{(ii) Parametric suppression from Planck operators.}

Using $\mu_{0i}^2 \sim \xi_i\, \Lambda_D^4/M_{\text{Pl}}^2$ and
$M_i^2 \sim \Lambda_D^2$:
\[
  \frac{\mu_{0i}^2}{M_i^2}
  \;\sim\; \frac{\xi_i\, \Lambda_D^4 / M_{\text{Pl}}^2}{\Lambda_D^2}
  \;=\; \xi_i\, \frac{\Lambda_D^2}{M_{\text{Pl}}^2}\,.
\]
Since $m_f = y_f\, v$ and $m_t = y\, v$:
\[
  \frac{m_f}{m_t}
  \;=\; \frac{y_f}{y}
  \;=\; \frac{\mu_{0i}^2}{M_i^2}
  \;\sim\; \xi_i\, \frac{\Lambda_D^2}{M_{\text{Pl}}^2}\,.
\]
\[
  \boxed{\frac{m_f}{m_t}
  \;\sim\; \xi_i\, \frac{\Lambda_D^2}{M_{\text{Pl}}^2}\,.}
\]

\textbf{Dimensional check:}
$[\Lambda_D^2/M_{\text{Pl}}^2] = [\text{dimensionless}]$ and
$[m_f/m_t] = [\text{dimensionless}]$.
\quad$\checkmark$

\textit{(iii) Numerical evaluation.}

\[
  \frac{\Lambda_D^2}{M_{\text{Pl}}^2}
  \;=\; \frac{(10^{11})^2}{(1.22 \times 10^{19})^2}
  \;=\; \frac{10^{22}}{1.49 \times 10^{38}}
  \;\approx\; 6.7 \times 10^{-17}\,.
\]

This is a \emph{single} insertion of the Planck suppression factor.
The observed mass ratios require different numbers of such insertions
(i.e., different ``tiers'' in the VEV hierarchy):

\renewcommand{\arraystretch}{1.3}
\begin{center}
\begin{tabular}{@{} l c c c @{}}
\toprule
\textbf{Fermion $f$} & $m_f / m_t$ (observed) & $\log_{10}(m_f/m_t)$
  & \textbf{Tier} \\
\midrule
Bottom ($b$) & $\approx 0.024$ & $-1.6$ & 1st \\
Charm ($c$) & $\approx 0.007$ & $-2.2$ & 1st--2nd \\
Strange ($s$) & $\approx 5 \times 10^{-4}$ & $-3.3$ & 2nd \\
\bottomrule
\end{tabular}
\end{center}

A single Planck suppression gives
$\log_{10}(\Lambda_D^2/M_{\text{Pl}}^2) \approx -16$, which is far
too small to reproduce $m_b/m_t \approx 0.024$ or even
$m_c/m_t \approx 0.007$.

\textbf{Key observation.}
The na\"ive one-step Planck suppression
$(\Lambda_D/M_{\text{Pl}})^2 \sim 10^{-16}$ is far too small to
reproduce $m_b/m_t \sim 0.024$.  This indicates that the bottom
quark mass does \emph{not} arise from the generic Planck-suppressed
operator mechanism of Eq.~\eqref{eq:mass-ratio-suppression}.
Instead, the scalar democracy framework provides a different
mechanism: \textbf{direct mixing} between the dominant Higgs doublet
$H_{33}^u$ and the subdominant doublet $H_{33}^d$ through the scalar
potential.  In this route, the suppression is controlled by
$\mathcal{O}(1)$ coefficients in the Planck operators and by the
mass splittings among the dormant Higgs doublets, rather than by
powers of $\Lambda_D/M_{\text{Pl}}$.  We derive this alternative
route in part~(c).

For the lighter fermions (charm, strange, and below), the generic
Planck suppression \emph{does} apply, through cascaded insertions of
$\Lambda_D^2/M_{\text{Pl}}^2$:
the parametric structure---a cascade of
progressively smaller VEVs---provides a natural mechanism for the
mass hierarchy, with the specific ratios set by $\mathcal{O}(1)$
coefficients $\xi_i$.

\bigskip

\textbf{(c) Bottom-to-top mass ratio.}

\textit{(i) Why the bottom mass requires direct $H^u$--$H^d$ mixing.}

If the bottom quark mass arose from the generic Planck-suppressed
mechanism of part~(b), we would have
$m_b/m_t \sim (\Lambda_D/M_{\text{Pl}})^2 \sim 10^{-16}$, which is
fourteen orders of magnitude too small.  This rules out the generic
tier structure for the bottom quark.

Instead, $H_{33}^d$ acquires its VEV through \textbf{direct mixing}
with $H_{33}^u$ in the scalar potential.  The Planck operators
of Eq.~\eqref{eq:planck-op} generate an off-diagonal mass-squared
$\mu_b^2$ between these two doublets, while the diagonal mass of
$H_{33}^d$ is $M_{d,33}^2$.  The induced VEV ratio is:
\[
  \frac{\langle H_{33}^d \rangle}{\langle H_{33}^u \rangle}
  \;\sim\; \frac{\mu_b^2}{M_{d,33}^2}
  \;\sim\; c_b \times \mathcal{O}(10^{-2})\,,
\]
where $c_b \sim \mathcal{O}(1)$ and the ratio
$\mu_b^2/M_{d,33}^2$ is determined by the specific pattern of
$\mathcal{O}(1)$ coefficients in the Planck-scale operators.
Crucially, neither $\mu_b^2$ nor $M_{d,33}^2$ is constrained to
be a power of $\Lambda_D/M_{\text{Pl}}$; their ratio is an
$\mathcal{O}(10^{-2})$ number set by the UV dynamics.

\textit{(ii) Lighter fermion tiers.}

The charm and strange masses arise from the same direct-mixing
mechanism but applied to different entries of the 18-Higgs scalar
potential:
\begin{itemize}
  \item \textbf{Charm} ($m_c/m_t \approx 0.007 \sim 10^{-2.2}$):
    The charm quark mass comes from the $(2,2)$ sector VEV
    $\langle H_{22}^u \rangle$, induced by direct mixing with the
    leading Higgs $H_{33}^u$.  Since
    $m_c/m_t \sim \mathcal{O}(10^{-2})$, the suppression ratio
    $\mu_c^2/M_{22}^2$ is of the same order as for the bottom quark.
    Both bottom and charm sit in the \textbf{first tier} of the VEV
    hierarchy.
  \item \textbf{Strange} ($m_s/m_t \approx 5 \times 10^{-4}
    \sim 10^{-3.3}$): The strange quark requires an additional
    suppression of $\sim 10^{-1}$ beyond the first tier.  This arises
    from a \textbf{cascaded} mixing: the $(2,2)$ or $(2,3)$ Higgs
    doublet mixes with the leading Higgs, and the strange-sector
    Higgs then mixes with this intermediate doublet.  Two steps of
    $\mathcal{O}(10^{-2})$ mixing give
    $\sim 10^{-4}$, consistent with $m_s/m_t$.  The strange quark
    sits in the \textbf{second tier}.
\end{itemize}
\[
  \boxed{\text{First tier (bottom, charm):}\;
  m_f/m_t \sim \mu_f^2/M_f^2 \sim 10^{-2};\quad
  \text{Second tier (strange):}\;
  m_s/m_t \sim (\mu^2/M^2)^2 \sim 10^{-4}\,.}
\]

\textit{(iii) VEV tiers.}

The VEV hierarchy can be organised as follows:
\begin{itemize}
  \item \textbf{Zeroth tier ($\langle H \rangle \sim v$):}
    $H_{33}^u$ (top quark).  This is the leading VEV that triggers
    electroweak symmetry breaking.
  \item \textbf{First tier ($\langle H \rangle \sim 10^{-2}\, v$):}
    $H_{33}^d$ (bottom quark), possibly $H_{22}^u$ (charm quark).
    These are induced by one step of Planck mixing, with the
    suppression factor controlled by $\mu_b^2/M_d^2 \sim 10^{-2}$.
  \item \textbf{Second tier ($\langle H \rangle \sim 10^{-4}\, v$):}
    Strange quark and lighter fermions.  These arise from two steps
    of mixing (Lecture~8, Eq.~(10)), giving a further suppression.
  \item \textbf{Third tier and beyond:}
    Up and down quarks, electron.  Progressive further suppression.
\end{itemize}

The crucial point: the \emph{parametric structure} of the hierarchy is
set by the scalar potential and Planck-scale operators.  The
\emph{specific values} of each tier's suppression factor depend on the
$\mathcal{O}(1)$ coefficients $c^{abcd}$, which are free parameters.

\bigskip

\textbf{(d) Discussion.}

\textit{(i) Prediction or consistency check?}

The fermion mass hierarchy is a \textbf{parametric consistency check},
not a prediction.  The free parameters are:
\begin{enumerate}
  \item $y$: the universal magnetic Yukawa coupling ($\sim 1$).
  \item $c^{abcd}$: the $\mathcal{O}(1)$ Planck-scale operator
    coefficients ($\sim 10^2$ independent parameters).
  \item $\Lambda_D$: the duality scale ($\sim 10^{11}\;\text{GeV}$).
  \item $M_{\text{Pl}}$: the Planck mass (measured:
    $1.22 \times 10^{19}\;\text{GeV}$).
  \item The scalar quartic couplings $\lambda_i$ and mass parameters
    $M_i^2$ of the 18-Higgs potential.
\end{enumerate}

The dualised SM has comparable freedom to the Standard Model's Yukawa
sector ($\sim 20$ free parameters for quark and lepton masses, mixing
angles, and CP phases).  It does \emph{not} predict specific mass
values.
\[
  \boxed{\text{Parametric consistency, not prediction.  The }
  \mathcal{O}(1) \text{ coefficients } c^{abcd}
  \text{ are free parameters.}}
\]

\textit{(ii) What the dualised SM explains.}

The Standard Model treats the Yukawa hierarchy as an unexplained input:
$y_t \approx 1$, $y_b \approx 0.024$, $y_e \approx 3 \times 10^{-6}$,
spanning five orders of magnitude with no structural reason for the
pattern.

The dualised SM, \emph{if the duality holds}, provides a
\textbf{structural explanation}:
\begin{itemize}
  \item The hierarchy arises from a \textbf{VEV hierarchy} among the 18
    Higgs doublets, not from a hierarchy in the Yukawa couplings (which
    are universal).
  \item The VEV hierarchy is \textbf{naturally generated} by
    Planck-scale operators that mix the leading Higgs with the dormant
    doublets.
  \item The \textbf{flavour puzzle} is \emph{translated} from the
    Yukawa sector to the scalar sector, where it has a gravitational
    origin.
\end{itemize}

\textit{(iii) Weakest assumption.}

The weakest assumption is the \textbf{non-SUSY duality itself}: the
claim that the electric $\SU{2}$ theory confines and produces the SM
as its magnetic dual.  If this duality does not hold, the entire
framework---universal Yukawa, 18 Higgs doublets, scalar
democracy---loses its theoretical basis.  The VEV hierarchy mechanism
is valid group theory conditional on the field content, but the field
content itself rests on the conjectured duality.

\end{proof}


\bigskip

\begin{center}
\rule{0.6\textwidth}{0.4pt}

\textit{End of Exercise Set 5.  This exercise develops the fermion mass
hierarchy from the scalar democracy mechanism of Lectures~7--8.  The
key result is that $m_t \sim y\, v$ with $y \sim \mathcal{O}(1)$ is a
parametric consistency check, and the lighter fermion masses arise from
a cascade of induced VEVs driven by Planck-scale operators.  The
$\mathcal{O}(1)$ coefficients $c^{abcd}$ are free parameters---the
framework provides the mechanism for the hierarchy but not the specific
mass values.}
\end{center}


%% ========================================================================
%%  BIBLIOGRAPHY
%% ========================================================================
\begin{thebibliography}{99}

\bibitem{CacciapagliaSannino2024}
  G.~Cacciapaglia and F.~Sannino,
  ``The Dualised Standard Model,''
  \texttt{arXiv:2407.17281} [hep-ph] (2024).

\bibitem{HillMachado2019}
  C.~T.~Hill, P.~A.~N.~Machado, A.~E.~Thomsen, and J.~Turner,
  ``Scalar democracy,''
  \textit{Phys.\ Rev.\ D} \textbf{100} (2019) 015015
  [\texttt{arXiv:1902.07214}].

\bibitem{Seiberg1995}
  N.~Seiberg,
  ``Electric--magnetic duality in supersymmetric non-Abelian gauge
  theories,''
  \textit{Nucl.\ Phys.} \textbf{B435} (1995) 129
  [\texttt{hep-th/9411149}].

\end{thebibliography}

\end{document}
